% Options for packages loaded elsewhere
\PassOptionsToPackage{unicode}{hyperref}
\PassOptionsToPackage{hyphens}{url}
\documentclass[
  ignorenonframetext,
]{beamer}
\newif\ifbibliography
\usepackage{pgfpages}
\setbeamertemplate{caption}[numbered]
\setbeamertemplate{caption label separator}{: }
\setbeamercolor{caption name}{fg=normal text.fg}
\beamertemplatenavigationsymbolsempty
% remove section numbering
\setbeamertemplate{part page}{
  \centering
  \begin{beamercolorbox}[sep=16pt,center]{part title}
    \usebeamerfont{part title}\insertpart\par
  \end{beamercolorbox}
}
\setbeamertemplate{section page}{
  \centering
  \begin{beamercolorbox}[sep=12pt,center]{section title}
    \usebeamerfont{section title}\insertsection\par
  \end{beamercolorbox}
}
\setbeamertemplate{subsection page}{
  \centering
  \begin{beamercolorbox}[sep=8pt,center]{subsection title}
    \usebeamerfont{subsection title}\insertsubsection\par
  \end{beamercolorbox}
}
% Prevent slide breaks in the middle of a paragraph
\widowpenalties 1 10000
\raggedbottom
\AtBeginPart{
  \frame{\partpage}
}
\AtBeginSection{
  \ifbibliography
  \else
    \frame{\sectionpage}
  \fi
}
\AtBeginSubsection{
  \frame{\subsectionpage}
}
\usepackage{iftex}
\ifPDFTeX
  \usepackage[T1]{fontenc}
  \usepackage[utf8]{inputenc}
  \usepackage{textcomp} % provide euro and other symbols
\else % if luatex or xetex
  \usepackage{unicode-math} % this also loads fontspec
  \defaultfontfeatures{Scale=MatchLowercase}
  \defaultfontfeatures[\rmfamily]{Ligatures=TeX,Scale=1}
\fi
\usepackage{lmodern}
\ifPDFTeX\else
  % xetex/luatex font selection
\fi
% Use upquote if available, for straight quotes in verbatim environments
\IfFileExists{upquote.sty}{\usepackage{upquote}}{}
\IfFileExists{microtype.sty}{% use microtype if available
  \usepackage[]{microtype}
  \UseMicrotypeSet[protrusion]{basicmath} % disable protrusion for tt fonts
}{}
\makeatletter
\@ifundefined{KOMAClassName}{% if non-KOMA class
  \IfFileExists{parskip.sty}{%
    \usepackage{parskip}
  }{% else
    \setlength{\parindent}{0pt}
    \setlength{\parskip}{6pt plus 2pt minus 1pt}}
}{% if KOMA class
  \KOMAoptions{parskip=half}}
\makeatother
\setlength{\emergencystretch}{3em} % prevent overfull lines
\providecommand{\tightlist}{%
  \setlength{\itemsep}{0pt}\setlength{\parskip}{0pt}}
% Auto-generated by infrastructure/rendering/_pdf_latex_helpers.py::extract_math_font_preamble.
% Minimal header passed to Pandoc via -H so Beamer slides pick up the same math font
% as the combined-PDF path. Adjust the manuscript's preamble.md to change the font globally.
\usepackage{unicode-math}
\setmathfont{latinmodern-math.otf}

% Slide layout defaults for warning-clean scientific decks.
\usepackage{etoolbox}
\IfFileExists{xurl.sty}{\usepackage{xurl}}{}
\IfFileExists{seqsplit.sty}{\usepackage{seqsplit}}{\newcommand{\seqsplit}[1]{#1}}
\protected\def\breakseq#1{\seqsplit{#1}}
\protected\def\breaktt#1{\begingroup\ttfamily\seqsplit{#1}\endgroup}
\setlength{\emergencystretch}{6em}
\tolerance=5000
\hbadness=10000
\hfuzz=1pt
\setlength{\tabcolsep}{2pt}
\AtBeginEnvironment{longtable}{\tiny\renewcommand{\arraystretch}{0.86}\setlength{\tabcolsep}{1pt}}
\AtBeginEnvironment{tabular}{\tiny\renewcommand{\arraystretch}{0.86}\setlength{\tabcolsep}{1pt}}

% Natbib and cross-reference fallbacks — slides don't load natbib
% or cleveref, but combined-PDF manuscript prose may emit these
% commands. The fallback renders citations as a bracketed key list
% and cross-references as detokenized labels so slides don't fail on
% undefined control sequences or raw underscores. \providecommand is
% a no-op if packages load later.
\providecommand{\citep}[1]{[#1]}
\providecommand{\citet}[1]{#1}
\providecommand{\citealp}[1]{#1}
\providecommand{\citeauthor}[1]{#1}
\providecommand{\citeyear}[1]{#1}
\providecommand{\cref}[1]{\texttt{\detokenize{#1}}}
\providecommand{\Cref}[1]{\texttt{\detokenize{#1}}}
\usepackage{bookmark}
\IfFileExists{xurl.sty}{\usepackage{xurl}}{} % add URL line breaks if available
\urlstyle{same}
% fallback for those not using the hyperref driver hyperxmp:
\makeatletter
\@ifundefined{xmpquote}{\newcommand{\xmpquote}[1]{#1}}{}
\makeatother
\hypersetup{
  hidelinks,
  pdfcreator={LaTeX via pandoc}}

\author{\texorpdfstring{}{}}
\date{}

\begin{document}

\section{Introduction}\label{sec:introduction}

\begin{frame}[fragile,allowframebreaks]{Introduction}
This \breaktt{template\_methods\_paper} serves as the
\textbf{methods-paper exemplar} for the
\href{https://github.com/docxology/template}{Research Project Template}
ecosystem: a manuscript whose subject is a methodology --- here, a
domain language for specifying controlled methods --- rather than
results produced by running one. The prose, the labelled figures, and
the compiled-plan table are produced through the same auditable custody
chain every exemplar in this template uses: tested functions in
\texttt{src/}, a thin analysis script, and generated-variable-injected,
multi-format rendering.
\end{frame}

\begin{frame}[allowframebreaks]{Why a methods paper needs its own DSL}
\protect\phantomsection\label{why-a-methods-paper-needs-its-own-dsl}
A methods section in ordinary prose is ambiguous by construction: ``add
10 mL of water, then mix'' admits multiple readings of order, units, and
what ``mix'' means operationally. BPL \breakseq{[@bpl2026]} makes the case for
laboratory protocols directly: free-text instructions admit multiple
interpretations, unit errors and reagent mismatches surface only at the
bench, and re-executing a protocol on a different operator or instrument
introduces silent variation. Fowler frames the general remedy as a
\textbf{domain-specific language} \breakseq{[@fowler2010dsl]}: a small,
purpose-built notation whose vocabulary is restricted exactly to the
concepts the domain needs, so that what can be written down is exactly
what is intended.
\end{frame}

\begin{frame}[fragile,allowframebreaks]{What this exemplar borrows from
BPL, and what it generalizes}
\protect\phantomsection\label{what-this-exemplar-borrows-from-bpl-and-what-it-generalizes}
BPL's architecture is a compiler pipeline --- parse, semantic check,
lower, schedule, execute, export --- over a biology-native type system
(units, dimensional analysis, MW-aware concentration), staged validation
gates, and deterministic compilation with a stable plan hash. Three
design choices carry over directly into \breaktt{src/methods\_dsl/}:

\begin{enumerate}
\tightlist
\item
  \textbf{Intent over instruction.} BPL users write high-level intents
  (\texttt{transfer}, \texttt{add\_reagent}, \texttt{incubate}); a
  compiler lowers them to backend-specific primitives.
  \breaktt{src/methods\_dsl/vocabulary.py}'s \texttt{StepKind} enum is
  the same idea, generalized:
  \texttt{TRANSFER}/\texttt{ADD}/\texttt{MIX} name \emph{what} happens,
  never \emph{how} a particular backend performs it.
\item
  \textbf{Dimensional safety.} BPL's type system catches
  \texttt{mL\ +\ g} at compile time, not at the bench.
  \breaktt{src/methods\_dsl/units.py} implements the same guarantee with
  a small \texttt{Dimension}/\texttt{Quantity} system rather than a full
  unit library.
\item
  \textbf{Deterministic compilation.} Same source, same options, same
  plan hash. \breaktt{src/methods\_dsl/compiler.py::compile\_method}
  reproduces this with a canonical-JSON SHA-256 hash over a
  Kahn's-algorithm \breakseq{[@kahn1962topological]} schedule.
\end{enumerate}

What this exemplar does \textbf{not} carry over is BPL's text grammar
and parser: a \texttt{Method} here is constructed directly as frozen
Python dataclasses (\breaktt{src/methods\_dsl/model.py}), not parsed from
\texttt{.bpl} source. This keeps the DSL's discipline in its typed,
validated shape rather than in new concrete syntax --- appropriate for a
template exemplar's scope --- while the controlled vocabulary,
dimensional safety, and deterministic compilation generalize unchanged
from wet-lab protocols to any controlled procedure, demonstrated in
\breakseq{[@sec:results]} by one wet-lab-flavored method and one
instrument-calibration method.
\end{frame}

\begin{frame}[fragile,allowframebreaks]{Template architecture context}
\protect\phantomsection\label{template-architecture-context}
The project sits on the repository's three pillars:

\begin{enumerate}
\tightlist
\item
  \textbf{\breaktt{src/methods\_dsl/} library}: pure, side-effect-free
  dataclasses and functions --- no plotting, no file I/O, and (with one
  declared logging exception) no \texttt{infrastructure} imports. This
  purity is what makes the library forkable and trivially testable.
\item
  \textbf{\texttt{tests/} framework}: a zero-mock suite that exercises
  every gate, the compiler, and the exporters against real
  \texttt{Method} fixtures covering both the success path and every
  gate-failure mode.
\item
  \textbf{\texttt{docs/} knowledge base}: the correspondence with BPL's
  pipeline, the testing philosophy, and the operational rules that
  govern agents editing this tree.
\end{enumerate}
\end{frame}

\begin{frame}[fragile,allowframebreaks]{The worked examples}
\protect\phantomsection\label{the-worked-examples}
We specify two methods with \breaktt{all\_example\_methods()}
(\breaktt{src/methods\_dsl/examples\_methods.py}):
\texttt{PBSPreparation}, an original --- not copied from BPL's shipped
examples --- manual bench preparation in BPL's own domain, and
\breaktt{SensorCalibrationSweep}, a non-biology controlled procedure
mixing automated measurement with a human sign-off step. The second
example exists specifically to demonstrate that the DSL's vocabulary
generalizes beyond wet-lab protocols, as \breakseq{[@sec:introduction]} claims.
\end{frame}

\begin{frame}[fragile,allowframebreaks]{Reader's guide to the
manuscript}
\protect\phantomsection\label{readers-guide-to-the-manuscript}
\begin{itemize}
\tightlist
\item
  \textbf{\breakseq{[@sec:methodology]}} ties each pipeline stage to its module
  in \breaktt{src/methods\_dsl/}.
\item
  \textbf{\breakseq{[@sec:results]}} is artifact-centric: every reported number
  names the function or report file that produced it.
\item
  \textbf{\breakseq{[@sec:experimental\_setup]}} lists the controlled
  vocabulary and software environment.
\item
  \textbf{\breakseq{[@sec:reproducibility]}} records the artifact inventory and
  the exact commands to regenerate everything.
\item
  \textbf{\breakseq{[@sec:scope]}} states scope and related work so the
  exemplar is not mistaken for a general-purpose protocol-execution
  system.
\end{itemize}
\end{frame}

\end{document}
